\input{preamble}

\section*{Moller scattering 3}

The Casimir trick uses matrix arithmetic to sum over spin states.
\begin{align*}
\sum_{abcd}\mathcal M_t\mathcal M_t^*&=\frac{e^4}{t^2}
\operatorname{Tr}\left(
(\slashed p_4+m)\gamma^\mu(\slashed p_2+m)\gamma^\nu
\right)
\operatorname{Tr}\left(
(\slashed p_3+m)\gamma_\mu(\slashed p_1+m)\gamma_\nu
\right)
\\
\sum_{abcd}\mathcal M_t\mathcal M_u^*&=\frac{e^4}{tu}
\operatorname{Tr}\left(
(\slashed p_4+m)\gamma^\mu(\slashed p_2+m)\gamma^\nu
(\slashed p_3+m)\gamma_\mu(\slashed p_1+m)\gamma_\nu
\right)
\\
\sum_{abcd}\mathcal M_u\mathcal M_u^*&=\frac{e^4}{u^2}
\operatorname{Tr}\left(
(\slashed p_3+m)\gamma^\mu(\slashed p_2+m)\gamma^\nu
\right)
\operatorname{Tr}\left(
(\slashed p_4+m)\gamma_\mu(\slashed p_1+m)\gamma_\nu
\right)
\end{align*}

Recall that slash matrices are formed as
\begin{equation*}
\slashed p=p^\mu g_{\mu\nu}\gamma^\nu
=p^0\gamma^0-p^1\gamma^1-p^2\gamma^2-p^3\gamma^3
\end{equation*}

\href{https://georgeweigt.github.io/examples/moller-scattering-3-demo.html}{Eigenmath script}

\end{document}
